% ============================================================
% Sección 8: Automatización de procesos de código
% ============================================================
\section{Automatizar procesos de código con agentes}

\begin{frame}{De autocompletado a agente de código}
  \begin{itemize}
    \item \textbf{Autocompletado} (Copilot clásico): sugiere la siguiente
      línea, el humano decide todo lo demás.
    \item \textbf{Agente de código} (Claude Code, Cursor Agent, Devin...):
      puede \emph{planear, escribir, ejecutar, ver el error y corregir} —
      todo en un loop, con supervisión humana en distintos grados.
    \item La diferencia clave: el agente tiene \textbf{herramientas} para
      leer/escribir archivos y \textbf{ejecutar} lo que escribe.
  \end{itemize}
\end{frame}

\begin{frame}{El loop generar-ejecutar-corregir}
  \centering
  \begin{tikzpicture}[node distance=1.2cm]
    \node[cajaLLM] (gen) {Generar/editar código};
    \node[cajaHerramienta, right=2cm of gen] (run) {Ejecutar\\(tests, linter, script)};
    \node[cajaEntorno, below=of run] (obs) {Observar salida\\(error, stack trace)};
    \node[cajaAccento, below=of gen] (fin) {¿Pasa? $\to$ Terminar};

    \draw[flecha] (gen) -- (run);
    \draw[flecha] (run) -- (obs);
    \draw[flecha] (obs) -- node[below, etiquetaFlecha]{si falla, corregir} (gen);
    \draw[flecha] (obs) -- node[right, etiquetaFlecha]{si pasa} (fin);
  \end{tikzpicture}
  \vfill
  \footnotesize Este es, literalmente, ReAct aplicado a ingeniería de software:
  \texttt{Thought} (qué cambiar) $\to$ \texttt{Action} (editar/ejecutar) $\to$
  \texttt{Observation} (resultado).
\end{frame}

\begin{frame}{Ejemplos de tareas automatizables}
  \small
  \begin{itemize}
    \item \textbf{Refactors mecánicos a escala}: renombrar una función y
      actualizar todos sus usos en un repo grande.
    \item \textbf{Corrección de bugs guiada por tests}: correr la suite,
      leer el fallo, editar, repetir hasta que pase.
    \item \textbf{Generación de código repetitivo}: boilerplate, migraciones,
      adaptadores entre APIs.
    \item \textbf{Revisión de código}: detectar bugs, sugerir simplificaciones
      antes de mergear.
    \item \textbf{Automatización de CI/CD}: diagnosticar por qué falló un
      pipeline y proponer el fix.
  \end{itemize}
\end{frame}

\begin{frame}{Qué necesita un agente para hacer esto bien}
  \begin{itemize}
    \item Acceso de \textbf{lectura/escritura} al sistema de archivos.
    \item Una forma de \textbf{ejecutar comandos} (shell, tests, build).
    \item \textbf{Feedback verificable}: tests, linters, tipos — señales
      objetivas de si el cambio funcionó (no solo ``se ve bien'').
    \item \textbf{Límites y supervisión}: permisos, sandboxing, confirmación
      humana en acciones riesgosas (borrar, hacer push, etc.).
  \end{itemize}
\end{frame}

\begin{frame}{Riesgos a tener en cuenta}
  \begin{alertblock}{No es magia, hay que vigilarlo}
    \begin{itemize}
      \item Puede ``arreglar'' el síntoma y no la causa si el feedback es débil.
      \item Puede ejecutar comandos destructivos si no hay barreras.
      \item Requiere buen criterio sobre \textbf{qué automatizar sin supervisión}
        y qué requiere aprobación humana explícita.
    \end{itemize}
  \end{alertblock}
\end{frame}
